\chapter{Mixed-Integer Linear Programming for Planning}
\label{ch:ch22}
% Function names for pseudocode are prefixed with "Milp" to stay unique
% across the book.
\SetKwFunction{MilpBranchAndBound}{BranchAndBound}
\SetKwFunction{MilpSolveLP}{SolveRelaxation}

\begin{objectives}
\begin{itemize}
  \item Write a planning question as a linear program (LP), explain why its feasible set is a polyhedron and why the optimum sits at a vertex, and say what integer and binary variables add: choices, either-or conditions, on/off switches and assignments.
  \item Derive the big-$M$ encoding of ``stay outside this rectangle'' and of ``stay at least $d_{\min}$ apart'', choose $M$ tightly, and count how many binaries a multi-vehicle problem needs.
  \item Assemble the complete mixed-integer linear program (MILP) for minimum-fuel and minimum-time multi-vehicle trajectories with double-integrator dynamics, and linearise 1-norm and $\infty$-norm objectives with auxiliary variables.
  \item Trace branch-and-bound by hand: LP relaxation, branching on a fractional binary, bounding with the incumbent, and why the tree can grow exponentially.
  \item Solve a two-vehicle instance with \code{scipy.optimize.milp}, check the solution, formulate MAPF as an integer program, and judge from a scaling experiment when exactness is worth its price.
\end{itemize}
\end{objectives}

% ---------------------------------------------------------------------
\section{Why this matters for a drone swarm}
\label{sec:ch22-motivation}

Three inspection drones must cross a courtyard in which a $2\times4$~m
building block stands in the way. They start at rest, must be at rest at
their goals six seconds later, may accelerate at no more than $2$~m/s$^2$,
must never come closer than one metre to each other, and the battery budget
asks for the plan with the smallest total velocity change. Every planner of
the earlier chapters answers a version of this question, and every one of
them compromises somewhere. \cbs (\cref{ch:ch09}) returns optimal paths, but
on a grid with unit moves, not trajectories with accelerations. The model
predictive controller of \cref{ch:ch21} handles the dynamics and the input
limits exactly, but a collision constraint is \emph{non-convex}: ``stay out
of the block'' is the union of four half-planes, and the quadratic program
of \cref{ch:ch21} can only keep a linearised version of it around the current
guess. The local planners of \cref{ch:ch12,ch:ch13,ch:ch14} look one step
ahead. None of them can tell you that no cheaper safe plan exists.

Mixed-integer linear programming (MILP)\index{MILP|see{mixed-integer linear
program}}\index{mixed-integer linear program} can. It keeps the linear
dynamics and the linear bounds of an LP and adds \emph{binary variables}
that switch constraints on and off: one binary per side of the block per time
step says which side the drone passes on, and one per axis direction per pair
of drones says on which side of each other they fly. The solver then searches
the combinations of these switches with a bound that lets it discard most
of them unseen, and it returns the global optimum together with a proof of
optimality. This is exactly the formulation with which Schouwenaars, De Moor,
Feron and How \cite{schouwenaars2001mixed} and Richards and How
\cite{richards2002aircraft} planned aircraft and spacecraft trajectories, and
it is the reference solution against which the hybrid planner of
\cref{ch:ch24} is benchmarked in \cref{ch:ch25}. The price is that, in the
worst case, the search is exponential in the number of binaries: MILP is the
tool for small, safety-critical or offline problems, for verifying that a
heuristic is close to optimal, and for scheduling take-off and landing slots,
not for replanning a swarm of fifty drones ten times a second. The training
plan lists it as a Medium-priority topic of Week~11 (\cref{ch:appA}), next to
MPC.

The chapter first recalls linear programs and what integer variables buy
(\cref{sec:ch22-intuition,sec:ch22-problem}), then derives the big-$M$
encodings of obstacle avoidance and separation and assembles the full
multi-vehicle MILP (\cref{sec:ch22-model}). \Cref{sec:ch22-algorithm}
explains branch-and-bound and traces it on a tiny instance,
\cref{sec:ch22-example} solves a two-vehicle instance with
\code{scipy.optimize.milp}, and \cref{sec:ch22-properties} proves what the
method guarantees and where it is exponential. \Cref{sec:ch22-variants}
formulates MAPF as an integer program and compares it with \cbs,
\cref{sec:ch22-cost} measures the cost of exactness, and
\cref{sec:ch22-implementation,sec:ch22-drone} cover the implementation, the
pitfalls and the place of MILP in the drone system.

% ---------------------------------------------------------------------
\section{The idea in plain words}
\label{sec:ch22-intuition}

A \textbf{linear program}\index{linear program} asks for the point that
minimises a linear function subject to linear inequalities. Each inequality
cuts the space with a half-plane, so the feasible set is a
\textbf{polyhedron}\index{polyhedron}: a convex region with flat sides and
sharp corners (\cref{fig:ch22-relaxation}). Sliding the level line of the
objective across the polyhedron shows that the best value is always reached
at a corner, a \emph{vertex}. The simplex method walks from vertex to
neighbouring vertex, always improving; interior-point methods travel through
the inside and converge to the same vertex. Both solve LPs with hundreds of
thousands of variables in seconds, and both are reviewed in
\cref{sec:ch22-lp}.

Now insist that some variables be integers, for instance $b \in \set{0,1}$
meaning ``the drone passes to the left'' or ``to the right'' of a block. The
feasible set is no longer a polyhedron but the set of lattice points inside
one, and it is not convex: the midpoint of two feasible points is usually
infeasible. Sliding the level line now stops at a lattice point that need not
be near the LP vertex. \Cref{fig:ch22-relaxation} shows both: the LP optimum
is at the fractional vertex $(2.5, 3.75)$, rounding it gives an infeasible
point, and the integer optimum $(3,3)$ lies on a different side of the
polyhedron. Dropping the integrality requirement is called the
\textbf{LP relaxation}\index{LP relaxation}; its value is a bound on the
integer optimum, and this bound is what makes an exact search affordable.

\begin{figure}[tb]
  \centering
  \inputfigure{ch22/relaxation}
  \caption{A linear program with two variables. The feasible polyhedron
  (blue) is cut by three half-planes; the level lines of the objective
  $x + 2y$ (dashed) show that the LP optimum is the vertex $(2.5, 3.75)$ with
  value $10$. If $x$ and $y$ must be integers, only the lattice points inside
  the polyhedron (dots) are feasible. Rounding the LP vertex to $(2,4)$ or
  $(3,4)$ leaves the polyhedron; the integer optimum is $(3,3)$ with value
  $9$, and the LP value $10$ is an upper bound on it.}
  \label{fig:ch22-relaxation}
\end{figure}

The second ingredient is the \textbf{big-$M$ trick}\index{big-M@big-$M$
trick}. A drone that stays out of a rectangular block must satisfy
\emph{at least one} of four inequalities: left of the block, right of it,
below it or above it. An ``or'' of linear inequalities is not linear. But if
we add to each inequality a term $M b_j$ with a large constant $M$ and a
binary $b_j$, then $b_j = 1$ relaxes the inequality so much that it holds
everywhere in the workspace (it is switched off), while $b_j = 0$ enforces it.
Demanding $\sum_j b_j \le 3$ keeps at least one inequality switched on.
The same trick encodes ``stay at least $d_{\min}$ apart'' for every pair of
drones. Everything else in the trajectory problem, the double-integrator
dynamics, the acceleration and speed limits, the start and goal conditions,
is already linear, and the objective can be made linear with auxiliary
variables. The result is one MILP, and a solver that combines LP relaxations
with a systematic search over the binaries, \textbf{branch-and-bound}, finds
its global optimum.

\begin{keyidea}
Binary variables turn ``either this or that'' into linear constraints: with
the big-$M$ trick, $b_j = 1$ switches an inequality off and $b_j = 0$ keeps
it on. Planning with dynamics, obstacles and separation then becomes one
mixed-integer linear program. Its LP relaxation is a bound; branch-and-bound
fixes fractional binaries one at a time and uses the bound to discard whole
subtrees, so it finds the exact optimum, at a cost that is exponential in the
number of binaries in the worst case.
\end{keyidea}

% ---------------------------------------------------------------------
\section{Problem statement and notation}
\label{sec:ch22-problem}

\subsection{Linear programs}
\label{sec:ch22-lp}

\begin{definition}[Linear program]
\label{def:ch22-lp}
A \textbf{linear program (LP)}\index{linear program!definition} in the
variables $\vect{z} \in \R^n$ is
\begin{equation}
  \min_{\vect{z}} \; \vect{c}\T \vect{z}
  \qquad \st \mat{A}_{\mathrm{eq}}\vect{z} = \vect{b}_{\mathrm{eq}}, \quad
  \mat{A}\vect{z} \le \vect{b}, \quad
  \vect{l} \le \vect{z} \le \vect{u},
  \label{eq:ch22-lp}
\end{equation}
with a linear objective $\vect{c}\T\vect{z}$, linear equality and inequality
constraints, and bounds $\vect{l} \le \vect{z} \le \vect{u}$ in which
entries may be $\pm\infty$. The \textbf{feasible set} is the polyhedron
$P = \set{\vect{z} : \text{all constraints of \cref{eq:ch22-lp} hold}}$.
The LP is \textbf{infeasible} if $P = \emptyset$, \textbf{unbounded} if the
objective has no lower bound on $P$, and otherwise it has an optimal value
$J^*$ attained at a vertex of $P$ (if $P$ has vertices).
\end{definition}

The geometry behind the last sentence is reviewed in \cref{ch:appB}: a
polyhedron is convex, a linear function on a convex polyhedron with vertices
takes its minimum at a vertex, and every vertex is determined by $n$ of the
constraints holding with equality. The \textbf{simplex method}\index{simplex
method} exploits this. It starts at a vertex, moves along an edge to a
neighbouring vertex with a smaller objective, and stops when no neighbour is
better; because the polyhedron is convex, a vertex without a better neighbour
is optimal. Its worst case is exponential in $n$, but on practical problems
it needs a number of steps that grows roughly linearly with the constraints.
\textbf{Interior-point methods}\index{interior-point method} follow a path
through the interior of $P$ towards the optimal face and take a number of
iterations that grows only logarithmically with the required accuracy; they
are polynomial in the worst case. Both are treated in
\textcite{nocedal2006numerical} and, together with the convex geometry, in
\textcite{boyd2004convex}. For this chapter it is enough to know that an LP
with $10^5$ variables is solved in well under a second, that a solver
reports infeasibility reliably, and that the LP solver is called thousands
of times inside a MILP solver. The HiGHS solver behind
\code{scipy.optimize.milp} uses a dual simplex method for these calls
\cite{huangfu2018parallelizing}.

\subsection{Integer and binary variables}
\label{sec:ch22-integer}

\begin{definition}[Mixed-integer linear program]
\label{def:ch22-milp}
A \textbf{mixed-integer linear program (MILP)}\index{mixed-integer linear
program!definition} is an LP \cref{eq:ch22-lp} together with an index set
$I \subseteq \set{1,\dots,n}$ of variables that must be integers,
$z_j \in \Z$ for $j \in I$. A variable with $z_j \in \set{0,1}$ is
\textbf{binary}\index{binary variable}. The \textbf{LP relaxation} of a
MILP is the LP obtained by dropping the integrality requirement. If all
variables are integers the program is an \textbf{integer program
(IP)}\index{integer program}.
\end{definition}

Integer variables buy four kinds of modelling power, listed in
\cref{tab:ch22-gadgets}. A binary can select one of several
\emph{choices}; it can enforce an \emph{either-or} between constraints with
the big-$M$ trick of the next section; it can act as an \emph{on/off switch}
that forces a continuous quantity to zero (a motor that is either off or runs
within its limits, a landing pad that is either free or used); and a matrix
of binaries with row and column sums encodes an \emph{assignment} of drones
to slots or goals. Logical relations between binaries are linear too:
``$z_1$ implies $z_2$'' is $z_1 \le z_2$, ``at most one of'' is
$\sum_j z_j \le 1$, and ``at least one of'' is $\sum_j z_j \ge 1$. Any
finite combinatorial decision can be written this way; the modelling craft,
surveyed by \textcite{vielma2015mixed}, lies in writing it so that the LP
relaxation is as tight as possible.

\begin{table}[tb]
  \centering\small
  \caption{What integer variables buy. Every gadget is linear in the
  continuous variables $x, u$ and the binaries $b, z$; $M$ is a constant
  large enough to make a switched-off inequality hold everywhere.}
  \label{tab:ch22-gadgets}
  \begin{tabular}{@{}lll@{}}
  \toprule
  Need & Encoding & Example in this book\\
  \midrule
  choice among $r$ options & $z_1 + \dots + z_r = 1$, $z_j \in \set{0,1}$ & arrival step of a drone\\
  either $f(x) \le 0$ or $g(x) \le 0$ & $f(x) \le M b$, $g(x) \le M(1-b)$ & pass left or right of a block\\
  on/off switch & $0 \le u \le u_{\max} z$ & pad in use or not\\
  assignment of $i$ to $j$ & $\sum_j z_{ij} = 1$, $\sum_i z_{ij} \le 1$ & drones to landing slots\\
  implication, at most one & $z_1 \le z_2$, \; $\sum_j z_j \le 1$ & one drone per vertex and time\\
  \bottomrule
  \end{tabular}
\end{table}

% ---------------------------------------------------------------------
\section{Modelling obstacles, separation and the trajectory problem}
\label{sec:ch22-model}

\subsection{Obstacle avoidance with the big-\texorpdfstring{$M$}{M} trick}
\label{sec:ch22-bigm}

Consider one drone with position $\pos_k = (x_k, y_k)$ at time step $k$ and
one rectangular obstacle $O = [x_{\min}, x_{\max}] \times [y_{\min},
y_{\max}]$, already inflated by the drone radius as in \cref{ch:ch02}, inside
a rectangular workspace $W = [X_{\min}, X_{\max}] \times [Y_{\min}, Y_{\max}]$
that bounds all positions. The drone is outside $O$ at step $k$ if and only
if at least one of the four inequalities
\begin{equation}
  x_k \le x_{\min}, \qquad x_k \ge x_{\max}, \qquad y_k \le y_{\min}, \qquad y_k \ge y_{\max}
  \label{eq:ch22-four-sides}
\end{equation}
holds (\cref{fig:ch22-bigm}). Introduce four binaries $b_{k,1}, \dots,
b_{k,4} \in \set{0,1}$ and a constant $M > 0$, and write
\begin{equation}
\begin{aligned}
  x_k &\le x_{\min} + M b_{k,1}, & x_k &\ge x_{\max} - M b_{k,2},\\
  y_k &\le y_{\min} + M b_{k,3}, & y_k &\ge y_{\max} - M b_{k,4},\\
  b_{k,1} + b_{k,2} + b_{k,3} + b_{k,4} &\le 3.
\end{aligned}
\label{eq:ch22-bigm}
\end{equation}
If $b_{k,j} = 0$ the $j$-th inequality of \cref{eq:ch22-bigm} is the $j$-th
inequality of \cref{eq:ch22-four-sides}: the constraint is \emph{on}. If
$b_{k,j} = 1$ it reads $x_k \le x_{\min} + M$, and if $M$ is large enough
this holds for every $x_k$ in the workspace: the constraint is \emph{off}.
The sum constraint forbids switching all four off, so at least one side
condition is enforced. The binaries are the drone's \emph{discrete}
decision, on which side of the block it is at step $k$; the solver, not the
modeller, makes it.

\begin{figure}[tb]
  \centering
  \inputfigure{ch22/bigm}
  \caption{The big-$M$ encoding of ``stay outside the rectangle $O$''. The
  drone at $\pos_k$ satisfies $x_k \le x_{\min}$, so $b_{k,1} = 0$ keeps that
  inequality on and the other three may be switched off ($b_{k,j} = 1$). A
  switched-off inequality, here $x_k \ge x_{\max} - M$ with the tight
  $M = x_{\max} - X_{\min}$, moves its boundary (dashed) to the edge of the
  workspace $W$ and is satisfied by every point of $W$. A larger $M$ moves
  the dashed line further out without gaining anything.}
  \label{fig:ch22-bigm}
\end{figure}

\paragraph{The meaning of $M$.}
For each inequality, $M$ must be at least the largest amount by which it can
be violated inside the workspace; this makes the switched-off inequality
vacuous and nothing more. The smallest such values are
\begin{equation}
  M_1 = X_{\max} - x_{\min}, \quad
  M_2 = x_{\max} - X_{\min}, \quad
  M_3 = Y_{\max} - y_{\min}, \quad
  M_4 = y_{\max} - Y_{\min},
  \label{eq:ch22-tightM}
\end{equation}
and \cref{prop:ch22-bigm} below shows that with these values
\cref{eq:ch22-bigm} is \emph{exactly} equivalent to ``outside $O$ at step
$k$''. Nothing forbids a larger $M$, and textbooks often write a single
huge constant, but a larger $M$ hurts twice.\index{big-M@big-$M$ trick!choice of M@choice of $M$}
First, it weakens the LP relaxation. In the relaxation the binaries are
fractional, and a point $x_k$ inside the obstacle satisfies
$x_k \le x_{\min} + M b_{k,1}$ with $b_{k,1} = (x_k - x_{\min})/M$; the larger
$M$, the closer this is to zero, so the relaxation can place the drone inside
the obstacle at almost no cost in the sum constraint and its bound tells
branch-and-bound nothing. Second, it damages the numerics. A solver accepts
a binary as integral when it is within a tolerance $\eps$ of $0$ or $1$,
typically $\eps = 10^{-6}$; a value $b_{k,1} = \eps$ is then reported as
$0$ although it relaxes the constraint by $M\eps$, which is one metre when
$M = 10^{6}$. Coefficients that differ by six orders of magnitude within one
row also make the simplex pivots ill-conditioned. Use \cref{eq:ch22-tightM},
computed from the actual workspace, and let the position bounds of the
workspace be explicit variable bounds so that the solver's presolve can
verify them.

\begin{proposition}[Exactness of the big-$M$ encoding]
\label{prop:ch22-bigm}
Let $\pos_k \in W$ and let $M_1, \dots, M_4$ be given by \cref{eq:ch22-tightM}
(or any larger values). Then $\pos_k \notin \operatorname{int} O$ if and
only if there exist $b_{k,1}, \dots, b_{k,4} \in \set{0,1}$ that satisfy
\cref{eq:ch22-bigm}.
\end{proposition}
\begin{proof}
If $\pos_k$ is not in the interior of $O$, one of the four inequalities
\cref{eq:ch22-four-sides} holds, say the $j$-th. Set $b_{k,j} = 0$ and the
other three binaries to $1$. The $j$-th row of \cref{eq:ch22-bigm} is the
$j$-th inequality, which holds; a row with $b = 1$ holds because $\pos_k \in
W$, for instance $x_k \le X_{\max} = x_{\min} + M_1$; and the sum is $3$.
Conversely, if binaries satisfying \cref{eq:ch22-bigm} exist, the sum
constraint forces some $b_{k,j} = 0$, whose row is the $j$-th inequality of
\cref{eq:ch22-four-sides}, so $\pos_k$ is not in the interior of $O$.
\end{proof}

Boundary points count as outside; if you need a strict margin, inflate the
obstacle. A convex polygon with $r$ edges is handled in the same way with
one binary per edge and $\sum_j b_j \le r - 1$, a box in three dimensions
with six binaries, and a moving obstacle with a predicted position at every
step (\cref{ch:ch20}) with a different rectangle $O_k$ per step. The
constraint is imposed at the sampled instants $k = 1, \dots, N$ only; what
happens between them is the subject of \cref{prop:ch22-intersample} and of
the pitfalls in \cref{sec:ch22-implementation}.

\subsection{Inter-vehicle separation}
\label{sec:ch22-separation}

Two drones $i$ and $j$ with positions $\pos^i_k$ and $\pos^j_k$ are separated
by at least $d_{\min}$ in the $\infty$-norm, $\norm{\pos^i_k -
\pos^j_k}_\infty \ge d_{\min}$, if and only if one of
\begin{equation}
  x^i_k - x^j_k \ge d_{\min}, \quad
  x^j_k - x^i_k \ge d_{\min}, \quad
  y^i_k - y^j_k \ge d_{\min}, \quad
  y^j_k - y^i_k \ge d_{\min}
  \label{eq:ch22-sep-sides}
\end{equation}
holds: drone $i$ is far enough to the right, left, above or below drone
$j$.\index{separation constraint!MILP} This is the obstacle condition again,
with a square of side $2 d_{\min}$ centred on the other drone, and it gets
the same treatment with four binaries $c^{ij}_{k,1}, \dots, c^{ij}_{k,4}$
per pair and step:
\begin{equation}
  x^i_k - x^j_k \ge d_{\min} - M c^{ij}_{k,1}, \quad
  x^j_k - x^i_k \ge d_{\min} - M c^{ij}_{k,2}, \quad
  y^i_k - y^j_k \ge d_{\min} - M c^{ij}_{k,3}, \quad
  y^j_k - y^i_k \ge d_{\min} - M c^{ij}_{k,4}, \quad
  \sum_{q=1}^{4} c^{ij}_{k,q} \le 3.
  \label{eq:ch22-sep}
\end{equation}
The tight constant is now $M = d_{\min} + (X_{\max} - X_{\min})$ for the
$x$-rows and $M = d_{\min} + (Y_{\max} - Y_{\min})$ for the $y$-rows, because
$x^i_k - x^j_k$ can be as small as $-(X_{\max} - X_{\min})$. The
$\infty$-norm square is a conservative stand-in for the Euclidean disc of
radius $d_{\min}$ (it contains the disc); the disc itself can be
approximated by a polygon with more binaries per pair.

\begin{table}[tb]
  \centering\small
  \caption{Number of binary variables $4N\bigl(m\,n_O + \binom{m}{2}\bigr)$
  for $m$ drones, one obstacle ($n_O = 1$) and $N$ time steps, and the
  number of leaves $2^{\text{binaries}}$ that a search without bounding
  would have to visit. The pair term $\binom{m}{2}$ dominates from about
  four drones on.}
  \label{tab:ch22-binaries}
  \begin{tabular}{@{}rrrrr@{}}
  \toprule
  Drones $m$ & Pairs $\binom{m}{2}$ & $N = 10$ & $N = 20$ & $N = 50$\\
  \midrule
  2 & 1 & 120 & 240 & 600\\
  3 & 3 & 240 & 480 & 1\,200\\
  5 & 10 & 600 & 1\,200 & 3\,000\\
  10 & 45 & 2\,200 & 4\,400 & 11\,000\\
  20 & 190 & 8\,400 & 16\,800 & 42\,000\\
  \bottomrule
  \end{tabular}
\end{table}

\Cref{tab:ch22-binaries} counts the binaries. With $m$ drones, $n_O$
rectangular obstacles and $N$ steps there are $4 N m n_O$ obstacle binaries
and $4 N \binom{m}{2}$ separation binaries; the pair term grows
quadratically in $m$, and every binary can in principle double the size of
the search tree. Two drones and twenty steps are a routine problem, ten
drones and fifty steps are a hard one, and twenty drones are beyond reach
without further structure. This table is the quantitative reason for the
verdict of \cref{sec:ch22-cost}.

\subsection{The full multi-vehicle trajectory problem}
\label{sec:ch22-formulation}

\begin{definition}[Multi-vehicle trajectory MILP]
\label{def:ch22-problem}
Given $m$ drones with start positions $\pos^i_{\mathrm{s}}$ and goals
$\pos^i_{\mathrm{g}}$, obstacles $O_1, \dots, O_{n_O}$ and a workspace $W$,
a horizon of $N$ steps of length $\dt$, acceleration and speed limits
$a_{\max}$ and $v_{\max}$ and a separation $d_{\min}$, find accelerations
$\ctrl^i_k \in \R^2$, $k = 0, \dots, N-1$, that minimise a cost $J$ subject to
\begin{subequations}
\label{eq:ch22-full}
\begin{align}
  \state^i_{k+1} &= \mat{A}\,\state^i_k + \mat{B}\,\ctrl^i_k,
    & &i = 1..m,\; k = 0..N-1, \label{eq:ch22-full-dyn}\\
  \state^i_0 &= (\pos^i_{\mathrm{s}}, \vect{0}), \quad
  \state^i_N = (\pos^i_{\mathrm{g}}, \vect{0}), & &i = 1..m, \label{eq:ch22-full-bc}\\
  \norm{\ctrl^i_k}_\infty &\le a_{\max}, \quad \norm{\vel^i_k}_\infty \le v_{\max},
  \quad \pos^i_k \in W, & &\text{all } i, k, \label{eq:ch22-full-bounds}\\
  \text{\cref{eq:ch22-bigm}} &\text{ for every } i, \text{ every } O_o
  \text{ and } k = 1..N, & &b^{io}_{k,j} \in \set{0,1}, \label{eq:ch22-full-obs}\\
  \text{\cref{eq:ch22-sep}} &\text{ for every pair } i < j \text{ and } k = 1..N,
  & &c^{ij}_{k,q} \in \set{0,1}, \label{eq:ch22-full-sep}
\end{align}
\end{subequations}
where $\state^i_k = (\pos^i_k, \vel^i_k) \in \R^4$ and $\mat{A}, \mat{B}$
are the zero-order-hold double-integrator matrices of \cref{ch:ch02},
\[
  \mat{A} = \begin{pmatrix} \mat{I} & \dt\,\mat{I}\\ \mat{0} & \mat{I}\end{pmatrix},
  \qquad
  \mat{B} = \begin{pmatrix} \tfrac12 \dt^2\,\mat{I}\\ \dt\,\mat{I}\end{pmatrix}.
\]
\end{definition}

Every constraint in \cref{eq:ch22-full} is linear: the dynamics are
equalities in the states and inputs, the norm bounds are boxes on single
components (the $\infty$-norm box is a linear inner approximation of the
Euclidean disc $\norm{\ctrl}_2 \le a_{\max}$ only if $a_{\max}$ is divided by
$\sqrt2$; the code keeps the box, as \textcite{richards2002aircraft} do), and
the obstacle and separation constraints are the big-$M$ rows. The states
could be eliminated by substituting the dynamics, but keeping them as
variables with sparse equality rows is what every modern solver prefers.

\paragraph{Linear objectives.}
A quadratic cost as in \cref{ch:ch21} is not allowed, but the two norms that
matter for a drone are. The \textbf{minimum-fuel}\index{minimum-fuel
objective} cost, the total velocity change $J_{\mathrm{fuel}} = \dt
\sum_{i,k} \norm{\ctrl^i_k}_1$, is linearised with one auxiliary variable
$s^i_{k,d} \ge 0$ per input component:
\begin{equation}
  J_{\mathrm{fuel}} = \dt \sum_{i,k,d} s^i_{k,d}, \qquad
  -s^i_{k,d} \le u^i_{k,d} \le s^i_{k,d}.
  \label{eq:ch22-fuel}
\end{equation}
At the optimum $s^i_{k,d} = \abs{u^i_{k,d}}$, because a larger $s$ would only
cost more. The \textbf{minimum-peak} cost $J_{\mathrm{peak}} = \max_{i,k}
\norm{\ctrl^i_k}_\infty$ uses a single auxiliary $s$ with $-s \le u^i_{k,d}
\le s$ for all $i, k, d$: the $\infty$-norm needs one variable where the
1-norm needs one per component. The \textbf{minimum-time}\index{minimum-time
objective} cost of \textcite{richards2002aircraft} needs binaries again:
$y^i_k \in \set{0,1}$ marks the arrival step of drone $i$,
\begin{equation}
  J_{\mathrm{time}} = \sum_{i} \sum_{k=1}^{N} k\,\dt\; y^i_k, \qquad
  \sum_{k=1}^{N} y^i_k = 1, \qquad
  \abs{p^i_{k,d} - p^i_{\mathrm{g},d}} \le M_{\mathrm{g}}
  \Bigl(1 - \sum_{k' \le k} y^i_{k'}\Bigr) \quad \text{for all } k, d,
  \label{eq:ch22-time}
\end{equation}
with $M_{\mathrm{g}}$ the workspace diameter: once the drone has arrived
($\sum_{k' \le k} y^i_{k'} = 1$) it must stay at the goal, before that the
row is switched off. A small multiple of $J_{\mathrm{fuel}}$ is added as a
tie-breaker among plans with the same arrival times. All three objectives
are implemented in \code{code/ch22\_milp.py}.

% ---------------------------------------------------------------------
\section{The algorithm: branch-and-bound}
\label{sec:ch22-algorithm}

\textbf{Branch-and-bound}\index{branch-and-bound} (Land and Doig
\cite{landdoig1960automatic}) is the search that every MILP solver runs at
its core. \Cref{alg:ch22-bnb} is its textbook form for a minimisation
problem $P$ with binaries.

\begin{algorithm}[htb]
\caption{Branch-and-bound for a MILP with binary variables}
\label{alg:ch22-bnb}
\KwIn{a MILP $P$ (\cref{def:ch22-milp}) with binaries indexed by $I$}
\KwOut{an optimal solution $\vect{z}^*$ or \emph{infeasible}}
\Fn{\MilpBranchAndBound{$P$}}{
  $J_{\mathrm{inc}} \gets \infty$; $\vect{z}^* \gets$ \KwNil
    \tcp*{incumbent: best integral solution so far}
  \Open $\gets \set{(P, -\infty)}$ \tcp*{nodes: a subproblem and the bound of its parent}
  \While{\Open $\neq \emptyset$}{
    $(Q, \beta) \gets$ node of \Open with the smallest $\beta$\label{alg:ch22-bnb:pop}\;
    \If{$\beta \ge J_{\mathrm{inc}}$}{\KwContinue \tcp*{pruned by the parent's bound}}
    $(\vect{z}, J_{\mathrm{LP}}) \gets$ \MilpSolveLP{$Q$}\label{alg:ch22-bnb:lp}\;
    \If{$Q$ is infeasible}{\KwContinue \tcp*{pruned by infeasibility}}
    \If{$J_{\mathrm{LP}} \ge J_{\mathrm{inc}}$}{\KwContinue \tcp*{pruned by bound}\label{alg:ch22-bnb:bound}}
    \If{$z_j \in \set{0,1}$ for all $j \in I$}{
      $J_{\mathrm{inc}} \gets J_{\mathrm{LP}}$; $\vect{z}^* \gets \vect{z}$
        \tcp*{integral: new incumbent}\label{alg:ch22-bnb:inc}
      \KwContinue\;
    }
    $j \gets$ a binary with fractional $z_j$, e.g.\ the one closest to $\tfrac12$\label{alg:ch22-bnb:pick}\;
    add $(Q \cup \set{z_j = 0}, J_{\mathrm{LP}})$ and $(Q \cup \set{z_j = 1}, J_{\mathrm{LP}})$ to \Open\label{alg:ch22-bnb:branch}\;
  }
  \KwRet{$\vect{z}^*$ (\emph{infeasible} if $\vect{z}^* = $ \KwNil)}\;
}
\end{algorithm}

The algorithm maintains an \textbf{incumbent}\index{incumbent}, the best
solution with integral binaries found so far, and a set of open subproblems.
Each subproblem is the original MILP with some binaries fixed to $0$ or $1$.
Line~\ref{alg:ch22-bnb:pop} selects the open node with the smallest parent
bound (\emph{best-first}); depth-first selection, which finds an incumbent
sooner, is the common alternative, and solvers mix both.
Line~\ref{alg:ch22-bnb:lp} solves the LP relaxation of the node, which is
cheap and gives $J_{\mathrm{LP}}$, a \emph{lower bound} on every solution of
the node with integral binaries (\cref{thm:ch22-bnb}). Three things can
happen. The relaxation is infeasible, so no integral solution exists in the
node either. Or the bound is not better than the incumbent
(line~\ref{alg:ch22-bnb:bound}): nothing in this node can beat what we have,
and the whole subtree is discarded unseen; this is the
\textbf{bounding}\index{branch-and-bound!bounding} that makes the method
work. Or the relaxation happens to be integral, and it is the best solution
of the node, so it becomes the incumbent (line~\ref{alg:ch22-bnb:inc}).
Otherwise some binary is fractional, and line~\ref{alg:ch22-bnb:branch}
\textbf{branches}\index{branch-and-bound!branching}: the node is replaced by
two children in which that binary is fixed to $0$ and to $1$. The union of
the children contains every integral solution of the parent, so nothing is
lost, and each child's relaxation is at least as constrained as the
parent's, so the bounds can only rise. Which binary to branch on
(line~\ref{alg:ch22-bnb:pick}) is a heuristic decision; ``most fractional''
is the simplest, and production solvers estimate which branching raises the
bound most.

\begin{example}[Branch-and-bound on a tiny instance]
\label{ex:ch22-tiny}
One drone must fly from $(0,0)$ to $(4,0)$, at rest at both ends, in $N = 4$
steps of $\dt = 1$ with $a_{\max} = 1$, $v_{\max} = 2$, the workspace
$[-1,5] \times [-3,3]$ and the inflated obstacle $O = [1.5, 2.5] \times
[-0.5, 0.5]$ on the straight line, with the minimum-fuel objective. The MILP
has $16$ binaries, four per step. Its LP relaxation flies straight through
the block with cost $4$ (the accelerations $(1,1,-1,-1)$ along $x$) and sets
every binary of step $k = 2$, where the drone is at $(2,0)$ inside $O$, to
$b_{2,j} = 0.143 = 0.5/3.5$: each side inequality is switched off by one
seventh, which the relaxation allows and the drone cannot do. The
implementation in \code{code/ch22\_milp.py} rounds the harmless binaries
first: in a group whose sum is at most $3$, a binary that is already $0$
certifies the disjunction, and the other three can be set to $1$ at no cost.
This leaves the group of step $2$, and \cref{tab:ch22-trace} and
\cref{fig:ch22-bnb} show the search. Fixing $b_{2,1} = 0$ demands
$x_2 \le 1.5$, which the dynamics cannot reach while still arriving at $x_4 =
4$ at rest, so that child is infeasible; $b_{2,2} = 0$ demands $x_2 \ge 2.5$,
also infeasible. The branch $b_{2,3} = 0$ (pass below, $y_2 \le -0.5$) has an
integral relaxation of cost $6$: the extra $2$ pays for the lateral
excursion $y = (0, -0.333, -0.5, -0.167, 0)$. Its sibling $b_{2,3} = 1$
forces $b_{2,4} = 0$ through the sum constraint, ``pass above'', a mirror
image with the same cost $6$, and it is pruned by the bound. Seven nodes and
six LPs settle the problem; enumerating the $4^4 = 256$ ways of choosing one
active side per step, which the self-test does as a check, gives the same
optimum $6$.
\end{example}

\begin{table}[tb]
  \centering\small
  \caption{Trace of \cref{alg:ch22-bnb} on \cref{ex:ch22-tiny} (best-first,
  branching on the most fractional binary of the smallest index). $J_{\mathrm{LP}}$
  is the value of the node's LP relaxation.}
  \label{tab:ch22-trace}
  \begin{tabular}{@{}rrlrll@{}}
  \toprule
  Node & Parent & Fixed & $J_{\mathrm{LP}}$ & Action & Incumbent\\
  \midrule
  1 & -- & -- & 4 & fractional $b_{2,1} = 0.143$: branch & $\infty$\\
  2 & 1 & $b_{2,1} = 0$ & infeasible & pruned & $\infty$\\
  3 & 1 & $b_{2,1} = 1$ & 4 & fractional $b_{2,2} = 0.143$: branch & $\infty$\\
  4 & 3 & $b_{2,2} = 0$ & infeasible & pruned & $\infty$\\
  5 & 3 & $b_{2,2} = 1$ & 4 & fractional $b_{2,3} = 0.143$: branch & $\infty$\\
  6 & 5 & $b_{2,3} = 0$ & 6 & integral: new incumbent & 6\\
  7 & 5 & $b_{2,3} = 1$ & 6 & $J_{\mathrm{LP}} \ge$ incumbent: pruned & 6\\
  \bottomrule
  \end{tabular}
\end{table}

\begin{figure}[tb]
  \centering
  \inputfigure{ch22/bnb}
  \caption{The branch-and-bound tree of \cref{ex:ch22-tiny}. Each node shows
  the value of its LP relaxation. The two left children are infeasible, node
  6 is integral and becomes the incumbent, and node 7, the mirror-image plan
  above the block, is discarded because its bound $6$ cannot beat the
  incumbent $6$. The root bound $4$ is the cost of flying straight through
  the block: the LP relaxation of a big-$M$ model is weak.}
  \label{fig:ch22-bnb}
\end{figure}

Two lessons of the example carry over to every MILP of this chapter. The
root relaxation of a big-$M$ model is nearly the unconstrained problem, so
the bound becomes useful only after several binaries are fixed, and
symmetric solutions (above/below, $i$ left of $j$ or $j$ left of $i$) create
subtrees that the bound can prune only after an incumbent of equal cost
exists. Production solvers such as HiGHS add \emph{presolve} (tightening
$M$ from the variable bounds, removing redundant rows), \emph{cutting
planes} (extra valid inequalities that cut off fractional vertices;
the combination is called \textbf{branch-and-cut}\index{branch-and-cut}),
and \emph{primal heuristics} that find incumbents early; see
\textcite{wolsey1998integer}. They change the constants, not the worst case
of \cref{thm:ch22-complexity}.

% ---------------------------------------------------------------------
\section{A worked example}
\label{sec:ch22-example}

\begin{example}[Two drones crossing past a block]
\label{ex:ch22-two}
Drone A flies from $(1, 4)$ to $(9, 4)$ and drone B from $(9, 4.5)$ to
$(1, 4.5)$ (metres), both at rest at both ends, in a workspace $[0,10] \times
[0,8]$. A building block occupies, after inflation by the drone radius, the
rectangle $O = [4,6] \times [1.5, 5.5]$, which cuts both straight lines. The
horizon is $N = 12$ steps of $\dt = 0.5$~s, $a_{\max} = 2$~m/s$^2$,
$v_{\max} = 3$~m/s and $d_{\min} = 1$~m. The passage above the block
($y \ge 5.5$) is closer to both drones than the passage below ($y \le 1.5$),
so both want to use it, head-on. The instance is built and solved by
\code{code/ch22\_milp.py} (\cref{sec:ch22-implementation}); every number
below is printed by its self-test.
\end{example}

\Cref{tab:ch22-sizes} lists the size of the MILP for the three objectives
and \cref{fig:ch22-example} shows the minimum-fuel and the minimum-time
trajectories. With the fuel objective there are $2 \cdot 4 \cdot 13 = 104$
state variables, $48$ inputs, $48$ slacks, $96$ obstacle binaries and $48$
separation binaries, $344$ variables in all, and $388$ constraint rows:
$8$ initial and $8$ terminal conditions, $96$ dynamics rows, $96$ slack rows,
$120$ obstacle rows ($2$ drones $\times$ $12$ steps $\times$ $5$ rows) and
$60$ separation rows. HiGHS solves it in about $1.8$~s after $366$
branch-and-bound nodes. The optimal fuel is $J_{\mathrm{fuel}} = 12.02$~m/s
of total velocity change. Both drones go over the top; A climbs to $y =
6.5$ while B stays on the edge $y = 5.5$, and at step $k = 6$ ($t = 3$~s)
they pass each other with $\norm{\pos^A_6 - \pos^B_6}_\infty = 1.000$: the
separation constraint is active, and the binary $c_{6,3} = 0$ (``A above B'')
is the solver's discrete decision. Raising $d_{\min}$ shrinks the feasible
set and can only raise the optimum: the self-test finds $J_{\mathrm{fuel}} =
10.77$, $12.02$, $12.56$ and $12.56$~m/s for $d_{\min} = 0.5$, $1$, $1.5$ and
$2$~m. The last two are equal, so an optimal plan for $1.5$~m already keeps
$2$~m at every sampled instant.

\begin{table}[tb]
  \centering\small
  \caption{Size of the MILP of \cref{ex:ch22-two} and the effort of
  \code{scipy.optimize.milp} (HiGHS 1.x, one core) for the three objectives.
  Times vary by a few tenths of a second between runs.}
  \label{tab:ch22-sizes}
  \begin{tabular}{@{}lrrrrrl@{}}
  \toprule
  Objective & Variables & Constraints & Binaries & B\&B nodes & Time [s] & Optimum\\
  \midrule
  minimum fuel & 344 & 388 & 144 & 366 & 1.8 & $12.02$~m/s\\
  minimum peak & 298 & 388 & 144 & 1 & 0.13 & $0.889$~m/s$^2$\\
  minimum time & 368 & 486 & 168 & 1 & 0.43 & $4.5 + 4.5$~s\\
  \bottomrule
  \end{tabular}
\end{table}

\begin{figure}[tb]
  \centering
  \inputfigure{ch22/example}
  \caption{Optimal trajectories of \cref{ex:ch22-two}, sampled every
  $\dt = 0.5$~s (dots), for the minimum-fuel objective (left) and the
  minimum-time objective (right); the gray rectangle is the inflated block.
  Left: A (blue) climbs over B (orange) and the two pass at $t = 3$~s with
  exactly $d_{\min} = 1$~m of $\infty$-norm separation. Right: both drones
  arrive after $4.5$~s at full acceleration, hug the top edge of the block,
  and \emph{swap} their $x$-positions between $t = 2$ and $t = 2.5$~s: every
  sampled separation is at least $1.25$~m, yet the continuous paths cross.
  The constraints hold only where they are imposed.}
  \label{fig:ch22-example}
\end{figure}

The minimum-peak objective gives a smooth plan with $\norm{\ctrl}_\infty \le
0.889$~m/s$^2$ whose relaxation is already integral: the block hardly matters
when the drones fly slowly. The minimum-time objective (right panel) brings
both drones home in $4.5$~s each, the discrete optimum of the arrival
binaries, at the price of $16.0$~m/s of velocity change and of the behaviour
described in the caption: at $t = 2$~s A is at $x = 4.375$ and B at $x =
5.625$, both on the edge $y = 5.5$; at $t = 2.5$~s they have exchanged these
positions. The sampled $\infty$-distance is $1.25$~m at both instants, so the
MILP is satisfied, but the continuous-time separation, evaluated by
\code{continuous\_check} in the code on the exact zero-order-hold motion,
drops to $0$: the drones fly through each other. The minimum-fuel plan is
better but not clean either: between samples its separation drops to
$0.856$~m and drone B cuts the top-left corner of the block by $0.146$~m
between $t = 3$ and $t = 3.5$~s. \Cref{prop:ch22-intersample} gives the
margins that make the sampled constraints imply the continuous ones, and
\cref{exr:ch22-margin} applies them to this instance.

